\documentclass[12pt,a4paper]{article}
\usepackage[french]{babel}
\usepackage[T1]{fontenc}
\usepackage{amsmath,amssymb,amsthm}
\usepackage{geometry}
\usepackage{fancyhdr}
\usepackage{listings}
\geometry{a4paper, margin=1in}

\title{Analyse Rigoureuse du Th\'eor\`eme du Tournesol d'Erd\H{o}s-Rado}
\author{Charles EDOU NZE\thanks{chercheur ind\'ependant}}
\date{\today}

\newtheorem{theorem}{Th\'eor\`eme}[section]
\newtheorem{lemma}[theorem]{Lemme}
\newtheorem{definition}[theorem]{D\'efinition}
\newtheorem{corollary}[theorem]{Corollaire}

\begin{document}
\maketitle
\thispagestyle{fancy}
\fancyhf{}
\renewcommand{\headrulewidth}{0pt}
\cfoot{\footnotesize Charles EDOU NZE, chercheur ind\'ependant}

\begin{abstract}
Nous pr\'esentons une investigation rigoureuse du th\'eor\`eme du tournesol d'Erd\H{o}s-Rado, un r\'esultat fondamental en combinatoire extr\'emale. Un tournesol \`a $w$ p\'etales est une collection de $w$ ensembles tels que chaque paire d'ensembles ait la m\^eme intersection. Nous formalisons le probl\`eme \`a travers des d\'efinitions axiomatiques strictes, analysons les structures combinatoires sous-jacentes et d\'etaillons les arguments inductifs classiques de majoration \'etablis par Erd\H{o}s et Rado. De plus, nous construisons une architecture explicite pour l'autoformalisation dans Lean 4 afin de garantir une v\'erification symbolique.
\end{abstract}

\section{D\'efinitions Axiomatiques et \'Enonc\'e du Probl\`eme}
\begin{definition}
Soit $\mathcal{F}$ une famille d'ensembles. Un sous-syst\`eme $\mathcal{S} \subseteq \mathcal{F}$ constitu\'e de $w$ ensembles distincts $S_1, S_2, \dots, S_w$ est appel\'e un \textbf{tournesol} (ou $\Delta$-syst\`eme) \`a $w$ p\'etales s'il existe un ensemble c{\oe}ur $C$ tel que pour tout $1 \le i < j \le w$:
\begin{equation}
S_i \cap S_j = C \label{eq:sunflower}
\end{equation}
Les ensembles $S_i \setminus C$ sont disjoints et non vides pour $i = 1, \dots, w$ (si $C \neq S_i$).
\end{definition}
\begin{definition}
Le th\'eor\`eme du tournesol d'Erd\H{o}s-Rado affirme qu'il existe une fonction $f(k, w)$ telle que toute famille $\mathcal{F}$ d'ensembles, chacun de cardinalit\'e d'au plus $k$, doit contenir un tournesol \`a $w$ p\'etales si $|\mathcal{F}| > f(k, w)$.
\end{definition}

\section{Litt\'erature Contextuelle}
Les travaux fondateurs sur les tournesols ont \'et\'e introduits par Erd\H{o}s et Rado en 1960. Ils ont d\'emontr\'e que $f(k, w) \le k!(w-1)^k$. La litt\'erature r\'ecente explore des lemmes de tournesol robustes et des am\'eliorations de ces bornes sup\'erieures, en s'appuyant toujours sur les propri\'et\'es d'intersection des syst\`emes d'ensembles. La preuve classique utilise l'induction sur la taille des ensembles dans la famille, une technique analogue \`a l'application du principe des tiroirs de Dirichlet sur des domaines de taille born\'ee.

\section{Lemmes et Preuves \'Etape par \'Etape}

\subsection{Cas de Base pour l'Induction}
\begin{lemma} \label{lem:base_case}
Pour $k=1$ et tout $w \ge 2$, $f(1, w) \le w - 1$.
\end{lemma}
\begin{proof}
Soit $\mathcal{F}$ une famille d'ensembles o\`u chaque ensemble a une cardinalit\'e d'au plus $1$. Ainsi, chaque ensemble dans $\mathcal{F}$ est soit un singleton, soit l'ensemble vide.
Supposons que $|\mathcal{F}| > w - 1$. Puisque tous les ensembles dans $\mathcal{F}$ sont distincts, au plus un ensemble peut \^etre l'ensemble vide, et les autres doivent \^etre des ensembles singletons distincts.
Deux ensembles singletons distincts (ou un singleton et l'ensemble vide) sont disjoints. Par cons\'equent, leur intersection est l'ensemble vide $C = \emptyset$.
Si nous s\'electionnons $w$ ensembles de $\mathcal{F}$ (ce qui est possible puisque $|\mathcal{F}| \ge w$), ils satisfont l'\'equation \eqref{eq:sunflower} avec $C = \emptyset$. Par cons\'equent, ils forment un tournesol \`a $w$ p\'etales.
\end{proof}

\subsection{\'Etape Inductive : L'Argument du C{\oe}ur}
\begin{lemma} \label{lem:inductive_step}
Si le th\'eor\`eme est vrai pour des ensembles de taille jusqu'\`a $k-1$, il est vrai pour des ensembles de taille jusqu'\`a $k$, sp\'ecifiquement born\'e par $f(k, w) \le k!(w-1)^k$.
\end{lemma}
\begin{proof}
Supposons, par induction, que $f(k-1, w) \le (k-1)!(w-1)^{k-1}$. Soit $\mathcal{F}$ une famille d'ensembles de taille au plus $k$, telle que $|\mathcal{F}| > k!(w-1)^k$.

Soit $\mathcal{D}$ une famille maximale d'ensembles mutuellement disjoints au sein de $\mathcal{F}$. Nous consid\'erons deux cas bas\'es sur la cardinalit\'e de $\mathcal{D}$.

Cas 1 : $|\mathcal{D}| \ge w$.
Par d\'efinition de $\mathcal{D}$, les ensembles dans $\mathcal{D}$ sont mutuellement disjoints. Par cons\'equent, pour tout couple d'ensembles $A, B \in \mathcal{D}$, $A \cap B = \emptyset$. Choisir $w$ ensembles quelconques de $\mathcal{D}$ donne un tournesol \`a $w$ p\'etales de c{\oe}ur $C = \emptyset$.

Cas 2 : $|\mathcal{D}| < w$.
Soit $X$ la r\'eunion de tous les ensembles dans $\mathcal{D}$. Puisque chaque ensemble dans $\mathcal{D}$ a une cardinalit\'e d'au plus $k$, le nombre total d'\'el\'ements dans $X$ est strictement born\'e :
\begin{align}
|X| &\le |\mathcal{D}| \cdot k \nonumber \\
|X| &\le (w - 1)k \label{eq:bound_x}
\end{align}

Puisque $\mathcal{D}$ est maximale, aucun ensemble dans $\mathcal{F}$ ne peut \^etre disjoint de $X$. Sinon, un tel ensemble pourrait \^etre ajout\'e \`a $\mathcal{D}$, contredisant sa maximalit\'e. Ainsi, chaque ensemble $S \in \mathcal{F}$ contient au moins un \'el\'ement $x \in X$.

Par le principe des tiroirs, il doit exister un \'el\'ement $x \in X$ qui est contenu dans de nombreux ensembles de $\mathcal{F}$. Sp\'ecifiquement, soit $\mathcal{F}_x$ la sous-famille des ensembles contenant $x$. Puisque $\mathcal{F} = \bigcup_{x \in X} \mathcal{F}_x$, nous avons :
\begin{align}
\max_{x \in X} |\mathcal{F}_x| &\ge \frac{|\mathcal{F}|}{|X|} \nonumber \\
&> \frac{k!(w-1)^k}{(w-1)k} \nonumber \\
&= (k-1)!(w-1)^{k-1}
\end{align}

Soit $x^*$ un \'el\'ement atteignant ce maximum, et soit $\mathcal{F}' = \{ S \setminus \{x^*\} \mid S \in \mathcal{F}_{x^*} \}$. La famille $\mathcal{F}'$ est constitu\'ee d'ensembles de cardinalit\'e au plus $k-1$. De plus, $|\mathcal{F}'| = |\mathcal{F}_{x^*}| > (k-1)!(w-1)^{k-1}$.

Par l'hypoth\`ese de r\'ecurrence, $\mathcal{F}'$ doit contenir un tournesol \`a $w$ p\'etales, disons $\{ S_1', S_2', \dots, S_w' \}$ de c{\oe}ur $C'$. En rajoutant l'\'el\'ement $x^*$ \`a chaque ensemble, on obtient $S_i = S_i' \cup \{x^*\}$ pour $i = 1, \dots, w$. Ces ensembles $S_i \in \mathcal{F}$ forment un tournesol de c{\oe}ur $C = C' \cup \{x^*\}$.
\end{proof}

\section{Architecture pour l'Autoformalisation}
La v\'erification formelle des lemmes susmentionn\'es peut \^etre impl\'ement\'ee dans Lean 4. Les types fondamentaux sont sp\'ecifi\'es comme suit :

\begin{lstlisting}[basicstyle=\ttfamily\small, breaklines=true]
import Mathlib.Data.Finset.Basic
import Mathlib.Data.Fintype.Basic

def IsSunflower {\alpha : Type*} (F : Finset (Finset \alpha)) (core : Finset \alpha) : Prop :=
  forall A B, A \in  F -> B \in  F -> A \neq  B -> A \cap  B = core

def HasSunflower {\alpha : Type*} (F : Finset (Finset \alpha)) (w : Nat) : Prop :=
  \exists  S : Finset (Finset \alpha), S \subseteq  F \land  S.card = w \land  \exists  core, IsSunflower S core

theorem erdos_rado_base_case {\alpha : Type*} (F : Finset (Finset \alpha)) (w : Nat) (hw : w >= 2)
  (hsize : \forall  A \in  F, A.card <= 1) (hcard : F.card > w - 1) :
  HasSunflower F w := by
  admit

theorem erdos_rado_induction {\alpha : Type*} (F : Finset (Finset \alpha)) (k w : Nat) (hw : w >= 2)
  (hsize : \forall  A \in  F, A.card <= k) (hcard : F.card > k.factorial * (w - 1)^k)
  (ih : \forall  G : Finset (Finset \alpha), (\forall  B \in  G, B.card <= k - 1) \rightarrow  G.card > (k - 1).factorial * (w - 1)^(k - 1) \rightarrow  HasSunflower G w) :
  HasSunflower F w := by
  admit
\end{lstlisting}

Cette structure fait efficacement le pont entre les arguments combinatoires informels et les exigences symboliques rigoureuses de Lean 4.

\end{document}
